\chapter{Research Context and Thesis Scope}

\section{Context}

This thesis is positioned in interactive autonomous systems, where agents must
make decisions under competing objectives and limited information. The central
technical lens is games of ordered preference (GOOPs), which encode strict
objective priorities instead of weighted trade-offs \cite{lee:2025,zanardi:2022}.

The practical motivation is to bridge three levels that are usually separated:

\begin{enumerate}
    \item formal game-theoretic planning,
    \item deployable robotics software,
    \item inverse inference of other agents' hidden preferences.
\end{enumerate}

\section{Thesis question}

The unifying question is:

\begin{quote}
How can ordered-preference game formulations be made practical for real
interactive systems when computation is constrained and other agents do not
reveal their preference structure?
\end{quote}

\section{Subproblems addressed}

This thesis decomposes that question into three subproblems.

\begin{description}
    \item[SP1: scalable forward solving] Efficient approximate solving of GOOPs
    in receding horizon.
    \item[SP2: embodied integration] Execution of the solver stack in
    Duckietown through a robust manager--solver--adapter architecture.
    \item[SP3: inverse preference estimation] Estimation of objective ordering
    from observed behavior in GOOP-like settings.
\end{description}

\section{Artifacts used in this thesis}

The manuscript is built around four concrete artifacts:

\begin{enumerate}
    \item the ITSC paper on approximate GOOP solving \cite{delasheras:2025},
    \item the \texttt{patos} integration repository,
    \item the \texttt{InverseGoopSolver} repository,
    \item the ongoing RA-L manuscript draft on preference estimation.
\end{enumerate}

\section{What this draft delivers}

At the current stage, this thesis draft provides:

\begin{itemize}
    \item full technical foundations and notation,
    \item complete software architecture description,
    \item implemented inverse-estimation workflow baseline,
    \item explicit experimental and writing roadmap to final submission.
\end{itemize}

Because parts of SP2 and SP3 are still in progress, this draft emphasizes
methodological clarity and reproducible structure rather than final
performance claims.

\begin{figure}[tbp]
    \centering
    \fbox{\parbox{0.92\linewidth}{
    \textbf{Placeholder figure: Thesis work map.}
    A map connecting SP1/SP2/SP3 with repositories, experiment scripts, and
    chapter outputs.
    }}
    \caption{Scope and structure of the MSc thesis work.}
    \label{fig:scope_map}
\end{figure}

\section{Contribution framing}

The contribution of this thesis is not a single isolated algorithm. Instead, it
is an integrated pipeline:

\begin{center}
\emph{approximate forward solving} $\rightarrow$
\emph{robotics integration} $\rightarrow$
\emph{inverse preference estimation}.
\end{center}

This framing is the key distinction from the earlier bachelor thesis and from a
paper-only manuscript.
